\section{Evaluation programme}

\subsection{Training arm}

The first phase uses small open models or controlled transformers so that the structural-redundancy mechanism can be tested before expensive scaling. Training datasets are generated or curated at controlled structural-redundancy levels while token length, difficulty and surface-semantic diversity are held as stable as feasible. Baselines include uniform sampling, exact/semantic deduplication or diversity selection, Skill-It-style dependency sampling, and MASS-style skill-graph selection. The RAKL arm differs only in using cross-domain structural objects and witnessed structural novelty rather than predefined surface skill labels.

Primary outcomes are tokens/examples to a frozen structural-OOD capability target, total training compute where measurable, performance on structure-known/domain-new tasks, performance on genuinely structure-new tasks, calibration and negative transfer. A positive result on ordinary in-domain accuracy without structural OOD does not support the central claim.

\subsection{Inference arm}

The main inference experiment freezes the base model. Baselines include vanilla reasoning, semantic exemplar/memory retrieval, analogy prompting, a reusable reasoning-primitive library and a structural workflow prior in the style of SWIFT. The RAKL arm retrieves a structural object, validates a directional witness, retrieves or adapts a validated operator where one exists, and verifies the target result. Outcomes include capability, Q2 transfer, Q3 invalid-transfer rate, reasoning tokens, latency, search nodes, tool calls and verification cost.

\subsection{Controlled structural redundancy}

A synthetic or semi-synthetic generator produces datasets with increasing effective duplication of underlying structures while surface forms remain varied. The preregistered curve is
\[
G(\rho_S)=\operatorname{CTC}_{\mathrm{parent}}-\operatorname{CTC}_{\mathrm{RAKL}},
\]
where $\rho_S$ is an uncertainty-aware structural-redundancy measure and CTC is cost to the frozen capability target. The analysis tests a predeclared monotone or rank relationship rather than choosing a curve after observing results.

\subsection{Parent-method assimilation}

Each strong baseline is first treated as a parent method. Its strongest relevant mechanism is reproduced as faithfully as feasible and mapped into the RAKL atlas. If the mechanism is useful, the final system may assimilate it; the scientific comparison then asks whether the residual RAKL structure adds value beyond the parent. For example, MASS and Skill-It become training-selection parents, SWIFT becomes a workflow-structure parent, and Reasoning Primitive Induction/TraceCompiler become procedural-reuse parents. The project does not intentionally weaken them to preserve novelty.

\subsection{Stopping rule}

Large-scale training is not authorized unless the low-cost mechanism tests pass. The programme stops or narrows if:
\begin{enumerate}[leftmargin=*]
\item the structural score does not add transfer information beyond semantic similarity;
\item Q3 semantic decoys are frequently transferred;
\item structure labels are unstable across independent annotators or extraction methods;
\item a parent method matches the cost/capability frontier;
\item induction and verification cost eliminate any realistic break-even point; or
\item the training and inference systems require unrelated representations.
\end{enumerate}

A null pilot is therefore useful: it prevents a large compute spend on a mechanism that has not survived its own easiest falsifiers.
